\chapter{What is actually measured}

This chapter states the pipeline from a deposited structure to a recovered edge, so that a reader can say what the number is a number about.

\section{From a deposited cell to exact points}

A CIF entry gives the cell lengths, the angles, and the fractional coordinates of the sites in one cell. Four steps turn that into something the detector reads.

\begin{enumerate}
\item \textbf{The angles are checked.} Only cells with right angles are used. An oblique cell would need a shear before a shift along one coordinate meant a shift along an axis, and that shear is a second thing able to go wrong between the published number and the recovered one.
\item \textbf{The cell is tiled.} A single cell holds one period and one period cannot be told from noise. Four tiles an axis, and there is no cap on this: the points are sparse. The cost is linear in the tile count and not the cube of a grid side.
\item \textbf{The sites are carried as exact integers.} A deposit writes decimal text, and that text is a numerator and a count of decimal places, both integers. A fractional coordinate plus a whole tile, times a cell edge, is a product of two decimals carried through as integers. Nothing rounds. There is neither grid, voxel nor nearest anything.
\item \textbf{Each axis is asked separately} for its period, and all three answers are kept. One structure therefore contributes up to three independent results.
\end{enumerate}

Nothing at any step tells the detector what it is looking at. It receives positions carrying values and an axis index.

\textbf{A grid is a cost the deposit does not impose.} Assigning sites to voxels of 0.25 angstroms and capping the tiling at 320 voxels along the longest side are both bounds from the reader and neither comes from the deposit. The evidence chapter measures what they cost: a mean absolute error of 0.0124 angstroms on a measurement whose error is otherwise zero. The cap also refuses large cells outright, a second cost that never appears as error because those entries go unread.

\section{Reading an axis without an order}

An axis is read as a lookup from a coordinate to the arrangement standing at it. Two coordinates agree when their arrangements match, values included, and a period is a coordinate whose arrangement matches the one exactly a period away.

The values have to be in the comparison. Leaving them out reads a rocksalt cell at half its published edge, correctly: the two sublattices interleave. The positions alone do repeat every $a/2$ and only the elements distinguish the two halves.

The candidates are every difference between two occupied coordinates, and that set is complete. A period that agrees with anything at all is by construction a difference between two things that agree. It is in the set. There is neither sweep, ceiling nor stride, and nothing in the measure narrows what can be found.

\textbf{Membership is the only relation asked of the set. The search is a lookup.} A binary search over sorted coordinates answers the same question and imposes an order the set does not have, paying a full comparison at every step of it.

\textbf{Taking the tallest lag is wrong.} A lattice of period $P$ agrees with itself just as well at $2P$ and at $3P$. Which of those is tallest is settled by noise. The tallest lag reports a harmonic and not the fundamental. Read that way, published cell edges come back at almost exactly twice their value on ten of eighteen axes.

The correction is to score a candidate on itself together with its multiples, and to require at least two members of that family to be present before the candidate is scored. Exactly two members are used and never more: a family growing as the candidate shrinks lets a short wrong candidate win by holding more members, since it only has to catch one good lag among them. An uncapped family reads herzenbergite, molybdite and one more mineral short by a factor near two thirds.

The family rule holds in the exact reading as it does on a grid. The grid's other costs are damage done by the reader. This one is a property of periodic sets themselves: a set with period $P$ agrees with itself at $2P$ whether anybody rounds it.

\textbf{This is the harmonic trap, and it is not specific to crystals.} Any reading of a period by its tallest peak falls into it. Scoring against multiples instead of taking the tallest peak is the general fix, and the crystal case is where it is easiest to see, because the right answer is published and a factor of two is unmistakable.

\section{What a failure would look like}

The criterion is stated before the numbers, since afterwards it is worth nothing.

A recovered period not equal to the published edge would mean the detector is not reading the period. A near miss counts as a miss, since nothing in this path rounds and there is no tolerance to fall inside. A systematic bias in one direction would mean it is reading the tiling and not the crystal, since the tiling is the only structure present that the mineral did not put there. Both are visible in the reported columns without knowing which entry is which.

A gridded reading has to set that bar at one voxel, because it has a voxel. It passes 453 of 453 and carries a bias of $-0.0067$ angstroms. The exact reading passes the stricter test and has no bias to explain.
